\subsection{Quadratic Fit Search}

\begin{frame}{Quadratic Fit Search: Explicación Gráfica}
    \begin{block}{Concepto principal}
        [Luna: Explica cómo se usan 3 puntos para trazar una parábola y cómo el vértice de esta aproxima el mínimo].
    \end{block}
    \begin{itemize}
        \item {[Idea gráfica 1]}
        \item {[Idea gráfica 2]}
    \end{itemize}
\end{frame}

\begin{frame}{Quadratic Fit Search: Pseudocódigo}
    \begin{algorithmic}[1]
        \State \textbf{Input:} [Entradas]
        \State [Fórmula para hallar el vértice de la parábola]
        \State [Lógica de reemplazo de puntos]
        \State \textbf{Return} [Resultado]
    \end{algorithmic}
\end{frame}

\begin{frame}{Quadratic Fit Search: Ejemplo Paso a Paso}
    \begin{exampleblock}{Planteamiento del problema}
        [Luna: Define aquí tu función y los 3 puntos iniciales de prueba].
    \end{exampleblock}
    \begin{itemize}
        \item \textbf{Paso 1:} [Cálculo del vértice]
        \item \textbf{Paso 2:} [Actualización de los 3 puntos]
    \end{itemize}
\end{frame}

\begin{frame}[fragile]{Quadratic Fit Search: Código}
\begin{lstlisting}[language=Python]
# [Luna: Reemplaza esto con tu código real]
def quadratic_fit_search_ejemplo():
    pass
\end{lstlisting}
\end{frame}

\begin{frame}{Quadratic Fit Search: Análisis de Convergencia}
    \begin{alertblock}{Tasa de Convergencia}
        [Luna: Menciona la convergencia superlineal de orden ~1.324].
    \end{alertblock}
    \begin{itemize}
        \item {[Condiciones en las que falla (ej. puntos colineales)]}
        \item {[Comparativa de eficiencia]}
    \end{itemize}
\end{frame}